Este capítulo apresenta uma avaliação honesta do estágio atual da pesquisa: o que foi concluído, o que permanece em aberto, quais desafios emergiram durante o desenvolvimento e o que se espera produzir na fase final. O objetivo não é apenas reportar progresso, mas identificar onde os riscos metodológicos ainda precisam ser mitigados antes da defesa.

\section{Progresso do Ciclo de Design Science Research}

Conforme a metodologia DSR adotada, o ciclo de pesquisa estrutura-se em seis fases. À data desta proposta (maio de 2026), quatro fases encontram-se concluídas ou em conclusão avançada:

\begin{enumerate}
    \item \textbf{Fase 1 - Identificação do Problema:} \textbf{Concluída}. Mapeamento das dificuldades de acesso aos dados parlamentares e definição da relevância para transparência e participação cívica foram estabelecidas.

    \item \textbf{Fase 2 - Objetivos da Solução:} \textbf{Concluída}. Especificação clara dos objetivos que o artefato deve atingir (redução de barreiras cognitivas, interface acessível, múltiplas modalidades) foi documentada.

    \item \textbf{Fase 3 - Design e Desenvolvimento:} \textbf{Concluída}. O artefato AgentAI encontra-se tecnicamente funcional, integrando coleta de dados (API Senado), processamento com IA Generativa (Mistral 7B local), e interfaces acessíveis (Dashboard + Chatbot).

    \item \textbf{Fase 4 - Demonstração:} \textbf{Concluída}. A aplicação encontra-se em ambiente de homologação, processando dados reais da API do Senado Federal e respondendo a consultas em linguagem natural.

    \item \textbf{Fase 5 - Avaliação:} \textbf{Em Execução}. Conforme cronograma (Capítulo \ref{cap:metodologia}, Quadro \ref{quad:cronograma_metodologia}), iniciaram-se preparações para testes técnicos, avaliação factual das respostas e verificação de rastreabilidade.

    \item \textbf{Fase 6 - Reflexão e Aprendizados:} \textbf{Planejada para Q3/Q4 2026}. Documentação do conhecimento científico gerado será realizada após conclusão dos testes.
\end{enumerate}

\section{Resultados Preliminares do Artefato}

Do ponto de vista técnico e funcional, o artefato demonstrou viabilidade nas seguintes dimensões:

\subsection{Viabilidade Técnica}

A solução computacional materializa-se como uma aplicação \textit{web} funcional que integra:

\begin{itemize}
    \item \textbf{Coleta de Dados:} Pipeline de ETL automatizado que processa dados da API do Senado Federal com sucesso, utilizando Python e \textit{Pandas} para estruturação.

    \item \textbf{Processamento com IA:} Modelo Mistral 7B Instruct, rodando localmente, consegue executar classificação temática de discursos legislativos e gerar respostas em linguagem natural alinhadas aos dados reais.

    \item \textbf{Interfaces:} Aplicação Streamlit com Dashboard interativo (\textit{Plotly}) e Chatbot conversacional funcionam de forma estável em ambiente de homologação.

    \item \textbf{Performance:} O sistema processa um corpus legislativo realista (sessões do Senado de 2024-2026) em ambiente local, com latência preliminarmente adequada ao estágio de prototipação.
\end{itemize}

\subsection{Verificação Funcional Preliminar}

Antes da fase de avaliação formal, o artefato passou por uma verificação funcional conduzida pela pesquisadora e pelo orientador, com foco exclusivo em identificar falhas de execução — não em medir desempenho. Nessa etapa, confirmou-se que: o pipeline de ETL processa e armazena dados da API sem perda de campos essenciais; o modelo Mistral 7B Instruct gera classificações temáticas para os discursos ingeridos; e a interface Streamlit responde a consultas sem erros de execução em ambiente local.

É importante distinguir essa verificação funcional da avaliação científica planejada para as próximas etapas. Saber que o sistema funciona não é o mesmo que saber se ele funciona bem o suficiente para o propósito declarado. A avaliação sistemática — com gabarito de referência anotado por dois avaliadores, métricas quantitativas e comparação com baseline — é o que determinará se a hipótese formulada se sustenta ou não.

\section{Desafios e Limitações Identificadas}

A execução das fases anteriores revelou desafios que constituem ameaças à validade da pesquisa e que serão tratados ou documentados nas próximas etapas:

\subsection{Desafios Técnicos}

\begin{itemize}
    \item \textbf{Ambiguidade Semântica em Jargões Legislativos:} O modelo Mistral 7B, apesar de competente em tarefas genéricas, demonstrou limitações na interpretação de termos específicos do domínio legislativo (ex: "desarquivamento", "tramitação em urgência", "voto por delegação"). Em alguns casos, isso resultou em classificações genéricas ou imprecisas.

    \textbf{Plano de Mitigação:} Engenharia de \textit{prompt} avançada com Few-Shot Learning (fornecer exemplos ao modelo de cada jargão) e possível fine-tuning limitado do modelo.

    \item \textbf{Alucinações e Opacidade dos LLMs:} A natureza "caixa preta" (\textit{black box}) de modelos de linguagem impõe um desafio para confiabilidade das respostas. Observou-se ocasionalmente que o Chatbot gera informações que parecem plausíveis mas não estão fundamentadas nos dados reais.

    \textbf{Plano de Mitigação:} Implementação de mecanismo de rastreabilidade: toda resposta do Chatbot cita quais dados (senador, data, votação) foram utilizados, permitindo verificação humana.

    \item \textbf{Limitações de Infraestrutura:} Restrições de memória e processamento em ambiente local impactam volume de dados que podem ser processados simultaneamente. Em alguns casos, há lentidão quando processando sessões legislativas muito longas.

    \textbf{Plano de Mitigação:} Otimização de código, caching de resultados, e possível upgrade de hardware.
\end{itemize}

\subsection{Desafios Metodológicos}

\begin{itemize}
    \item \textbf{Construção de Gabaritos de Referência:} Avaliar classificação temática e respostas em linguagem natural exige a construção de conjuntos de referência consistentes. Há risco de subjetividade na definição das categorias e respostas esperadas.

    \textbf{Plano de Mitigação:} Documentação explícita dos critérios de anotação, revisão pelo orientador e registro dos casos ambíguos encontrados durante a avaliação.

    \item \textbf{Validade dos Cenários de Consulta:} Os cenários usados para avaliar o agente conversacional precisam representar perguntas relevantes sobre dados parlamentares, sem depender de participantes externos.

    \textbf{Plano de Mitigação:} Definição de cenários com base nas funcionalidades reais do artefato, nos dados disponíveis na API do Senado e nas lacunas identificadas no mapeamento sistemático.
\end{itemize}

\section{Contribuições Esperadas da Pesquisa}

Conforme apontado pelo Mapeamento Sistemático da Literatura (Capítulo 4), esta pesquisa posiciona-se para fazer contribuições em áreas onde há lacunas identificadas:

\subsection{Contribuições Científicas}

\begin{enumerate}
    \item \textbf{Conhecimento sobre Design de Interfaces para Acessibilidade de Dados Legislativos:} Investigação fundamentada de como a combinação de múltiplas modalidades (Dashboard visual + Chatbot conversacional) pode favorecer as pré-condições técnicas de acesso à informação legislativa, reconhecendo-se que a verificação do efeito junto a cidadãos reais constitui trabalho futuro.

    \item \textbf{Compreensão de Efetividade Técnica de IA Generativa em Contextos Cívicos:} Avaliação rigorosa de como IA Generativa local pode apoiar classificação temática, consulta em linguagem natural e geração de respostas rastreáveis sobre dados parlamentares.

    \item \textbf{Metodologia para Validação de Confiabilidade de LLMs em Contextos Críticos:} Documentação de estratégia de validação de respostas IA contra dados reais, aplicável a outros domínios além legislação.

    \item \textbf{Conhecimento Localizado (Brasil/Português):} Investigação situada no contexto brasileiro, em língua portuguesa, contribuindo para uma lacuna ainda pouco explorada na literatura sobre IA Generativa e dados parlamentares.
\end{enumerate}

\subsection{Contribuições Práticas}

\begin{enumerate}
    \item \textbf{Artefato Funcional:} Protótipo funcional que pode subsidiar futuras aplicações por pesquisadores, organizações civis e instituições interessadas em transparência legislativa.

    \item \textbf{Código Aberto:} Publicação do código-fonte no GitHub permite que outros pesquisadores repliquem, modifiquem e estendam a solução.

    \item \textbf{Guidelines de Design:} Recomendações práticas para design de interfaces cívicas que combinam IA com acessibilidade.
\end{enumerate}

\section{Trabalhos Futuros: Fases de Avaliação e Conclusão}

As fases restantes do ciclo DSR — Avaliação e Reflexão — constituem o núcleo científico da dissertação. É nelas que a hipótese será testada e o conhecimento gerado. O cronograma detalhado das atividades está descrito no Capítulo \ref{cap:metodologia}, Quadro \ref{quad:cronograma_metodologia}; registra-se aqui apenas a priorização das atividades críticas:

A construção do protocolo de anotação e do gabarito de referência precede qualquer execução de experimento — sem um gabarito metodologicamente sólido, as métricas de acurácia e F1 carecem de validade. Em seguida, a execução dos testes de classificação com baseline comparativo permitirá dimensionar o real ganho do uso do LLM. Por fim, a avaliação de rastreabilidade e os cenários de consulta fecharão o ciclo de evidências sobre a hipótese central.

\section{Perspectivas de Conclusão}

A fase que determina o valor científico desta dissertação ainda está por ser executada: a avaliação rigorosa do artefato. O desenvolvimento do AgentAI demonstrou viabilidade técnica, mas viabilidade não é o mesmo que efetividade — e é exatamente a distinção entre essas duas condições que a fase de avaliação precisa endereçar.

Cabe reiterar o alcance do que se promete. O título e a motivação evocam a \textit{democratização do acesso}, mas o que esta dissertação efetivamente projeta e avalia são as \textbf{pré-condições técnicas} dessa democratização --- a qualidade factual, a rastreabilidade das respostas e a disponibilização de interfaces de baixa barreira de entrada. A demonstração do efeito democratizante junto a cidadãos reais, que exigiria um estudo com participantes, é reconhecida como trabalho futuro e não é reivindicada como resultado deste trabalho.

O principal risco no horizonte imediato é a construção do gabarito de referência: se os critérios de anotação não forem documentados com precisão antes do início da classificação, a avaliação perde credibilidade metodológica. Por isso, a primeira atividade do próximo trimestre é a elaboração do protocolo de anotação — incluindo definição operacional das categorias temáticas, exemplos de casos limítrofes e regras de desempate — antes que qualquer discurso seja anotado.

Se as fases de Avaliação e Reflexão forem executadas com o rigor descrito neste documento, o trabalho poderá contribuir com dois resultados concretos para a comunidade: um protocolo replicável para avaliação de sistemas RAG sobre dados parlamentares em português, e um conjunto de dados anotados sobre discursos do Senado Federal que pode ser disponibilizado para pesquisas futuras. Esses resultados têm valor independente do desempenho absoluto do modelo — porque documentam tanto o que funciona quanto o que não funciona, e ambos são contribuições científicas legítimas.

\textbf{Próximas Ações Imediatas:}

\begin{enumerate}
    \item Refinamento final do documento de proposta conforme feedback da banca examinadora
    \item Preparação do dataset de controle (Gabarito Ouro) para experimento de acurácia
    \item Preparação de cenários representativos de consulta
    \item Implementação de melhorias metodológicas apontadas por este parecer
    \item Execução da primeira rodada de avaliação técnica (próximas 8 semanas)
\end{enumerate}
